\section{课程定位：从模型能力竞争转向系统能力竞争}

\subsection{主讲人视角与问题定义}
这讲并不是在讲某个单点模型技巧，而是从系统工程历史去解释一个更大的命题：当大模型和 Agent 成为主流工作负载时，整个基础设施设计为什么必须重写。Yangqing Jia 的经历横跨 Google 研究、Caffe/TensorFlow/PyTorch 生态、Lepton AI 创业与被并入 NVIDIA 的工程实践，因此他的观察并非学术抽象，而是来自大规模生产系统中的成本、可靠性与演进压力。

\begin{importantbox}{本讲的核心问题}
我们今天讨论的不是 ``Agent 是不是新概念''，而是 ``当 Agent 工作负载进入真实生产环境，哪些系统设计假设会失效，哪些旧原则会回归，哪些新抽象必须补上''。
\end{importantbox}

课程从三个层次递进：
\begin{enumerate}
    \item 技术浪潮层：MLOps $\rightarrow$ RAG $\rightarrow$ Agentic AI 的热度切换与真实留存；
    \item 基础设施层：HPC、Cloud、Data Cloud、AI Cloud 的架构约束如何变化；
    \item 产业落地层：企业应用、创业机会、边缘部署、成本与价值之间如何权衡。
\end{enumerate}

\begin{knowledgebox}{为什么 ``系统设计'' 在 Agent 时代更关键}
传统 LLM 应用可以停留在 ``单轮请求-响应''。Agent 系统则包含任务拆解、状态管理、外部工具调用、长链路执行、失败恢复与审计闭环。模型只是其中一个部件，真正决定能否商用的是系统能否稳定承载这些环节。
\end{knowledgebox}

\begin{warningbox}{常见误区：把 Agent 仅理解为 Prompt 工程升级}
如果只在提示词层面做 ``更复杂的模板''，没有补齐调度、存储、检索、流控、可观测性与故障恢复，那么系统会在高并发和长任务下快速失稳。很多所谓 ``模型不够强'' 的问题，实际是系统设计不足导致的。
\end{warningbox}

\subsection{本章小结}
本讲的价值在于把 Agent 讨论从概念热度拉回工程现实：模型能力提升只是起点，系统演进才是规模化落地的决定因素。

